% GENERATED by tools/emit_compile_appendix.py -- do not edit by hand.
% Every number is read from data/cryptobench_apo/TABLE_FIELD.json and the
% dataset manifests. Regenerate with `make compileapp`.

\section{Appendix C: the compile, the boundary cases, and the units}
\label{app:compile}

Appendix~\ref{app:wires} states what the \(43\) local quantities are and how they expand into the \NTabWires{} wires. This states everything between a wire and a score, and then the two things a set of formulas cannot state: what happens at the boundaries, and which structures the numbers were computed over.

\subsection{The table bank and its seed}

The bank is 16 independent random partitions of the \NTabWires{} wires into groups of 2, drawn from \texttt{numpy.random.default\_rng} seeded with \(20260725\). The seed is recorded in the compiled artifact and the generator of this appendix redraws the bank from it and refuses to emit unless the redraw matches the shipped bank pair for pair.

A partition rather than a draw of random tuples: within one round every wire appears at most once, so no wire is over-represented, and across rounds each wire meets a different set of partners. Coverage of the digit-pair lattice comes from the number of rounds and not from the width of any table.

\NTabWires{} is odd, so a partition into pairs leaves one wire over in every round, and a group of one addresses no pair and is dropped. Each round therefore contributes 322 tables and the bank holds \(322 \times 16 = \NTabTables{}\). The consequence is worth stating because it is easy to assume otherwise: 629 wires appear in 16 tables and 16 appear in 15, one for each round, and which wire is left out is a property of the seed. Every wire is in the bank; not every wire is in it equally often.

\subsection{What a cell holds}

Each table \(k\) of width 2 over \(L = 4\) levels has \(L^{2} = 16\) addresses, for \NTabCells{} cells across the bank. A cell holds two integers counted on the training fold: \(\mathrm{tot}_k[a]\), the residues that addressed it, and \(\mathrm{pos}_k[a]\), those of them that are labelled cryptic binding residues. The value read at inference is
\begin{equation}
  \hat p_k[a] = \begin{cases}
    \mathrm{pos}_k[a] / \mathrm{tot}_k[a], & \mathrm{tot}_k[a] > 0,\\[2pt]
    \bar y = 0.057589, & \mathrm{tot}_k[a] = 0,
  \end{cases}
\end{equation}
where \(\bar y\) is the training fold's base rate. An address no training residue reached takes that base rate because it is the only value that adds nothing to the score, and there are \NTabCellsEmpty{} such cells. No smoothing is applied to the cells that were reached: a cell visited once holds \(0\) or \(1\), and the integer fan-out is what decides how much such a cell is allowed to say.

\subsection{The integer fan-out}

Write \(v_k(i) = \hat p_k[a_k(i)]\) for the value table \(k\) returns on residue \(i\), and \(v(i) \in \mathbb{R}^{\NTabTables{}}\) for the stacked vector. Let \(\mu_1\) and \(\mu_0\) be its means over the labelled and unlabelled training residues and \(S\) its pooled scatter. The fan-out is the rounded, rescaled solution of one regularised symmetric system:
\begin{align}
  \tilde S &= S + \left( \lambda \, \frac{\operatorname{tr} S}{K} + \varepsilon \right) I, \qquad \lambda = \TabRidge{}, \quad \varepsilon = 10^{-12}, \quad K = \NTabTables{},\\
  w &= \tilde S^{-1} (\mu_1 - \mu_0),\\
  m_k &= \operatorname{round}\!\left( \frac{w_k}{\max_j |w_j|} \, C \right), \qquad C = \TabCap{}.
\end{align}

Three points. The ridge is proportional to the mean diagonal of the scatter rather than absolute, so it means the same thing whatever scale the cell rates happen to sit on, and it is not cosmetic: pairs drawn from the lattice repeat as the pool grows, the scatter goes near-singular, and without regularisation the direction chases the null space. The normalisation is by the largest coordinate, not the norm, so the cap is attained by exactly the tables the solve ranked highest and the rounding grid is the same regardless of how many tables there are. And the rounding is to nearest, which sends every table whose weight is below half a grid step to zero --- which is why \NTabUsedFullFold{} of \NTabTables{} tables carry a non-zero multiplicity and the rest are compiled but never consulted. The solve is closed form: no gradient, no iteration, and nothing selected on the held-out fold.

\subsection{The gate and the operating point}

The raw score is \(S(i) = \sum_k m_k v_k(i)\). The gate adds the mean of \(S\) over the residues within \(18\)~\AA\ of residue \(i\)'s centroid, inside its own chain, rescaled so that the added term has the same standard deviation as \(S\) before it is weighted:
\begin{equation}
  F(i) = S(i) + 1 \cdot \bar S_{18}(i) \cdot \frac{\operatorname{sd}(S)}{\operatorname{sd}(\bar S_{18})}.
\end{equation}
Matching the standard deviation rather than the maximum matters because the maximum of a score field over a chain is an order statistic of a handful of residues, so a max-normalised gate would mix in a different amount on every structure. Both constants were chosen on a cluster-disjoint half of the training fold. If \(\operatorname{sd}(\bar S)\) is zero the gate is skipped, which can only happen on a chain whose residues all score alike.

The positive call is the top \(9\%\) of residues by \(F\) within each chain, \(k = \max(1, \operatorname{round}(q n))\), ties broken by ascending residue index so the rule is deterministic. The threshold \(q = 0.09\) is the argmax of the pooled F1 curve on the training fold and was frozen before the held-out fold was binarised.

\subsection{Boundary cases}

These are the cases the formulas above do not determine. Each is a decision in the code, and each is stated here because the alternative decisions are all defensible and a reader cannot tell which was taken.

\begin{description}
  \item[The neighbourhood is never empty] \(N_r(i)\) is \(\{j : \|c_j - c_i\| \le r\}\), which contains \(i\) itself, so no statistic ever divides by zero and there is no isolated-residue case to define. A residue with no other residue inside \(r\) therefore takes its own value as the mean, \(0\) for the dispersion, \(0\) for the centred difference and \(0\) for the local rank. Self-inclusion also means the local rank of the largest residue in its neighbourhood is \((|N_r(i)|-1)/|N_r(i)|\) rather than \(1\), and that the dispersion is a population standard deviation over a set that includes the point it is centred on.
  \item[Zero variance] Dispersions are computed as \(\sqrt{\max(\langle x^2 \rangle - \langle x \rangle^2, 0)}\), so floating-point cancellation on a neighbourhood of near-identical values cannot produce a negative variance or a NaN.
  \item[Ties in a rank] In the banding, tied values share the mid-rank of the run, so the digit assignment does not depend on the order the residues happen to appear in the file. A wire constant across a chain gives every residue the same mid-rank and lands them all in one digit, which addresses one cell per table and contributes the same amount to every residue's score --- visible in the ranking not at all, which is why the audit of Section~\ref{sec:audit} decomposes deviations from the chain mean rather than absolute scores.
  \item[Missing atoms] The residue universe is the distinct residue numbers among \texttt{ATOM} records whose element is not hydrogen, so a residue enters if it has at least one non-hydrogen atom, and its centroid is the centroid of the atoms actually present. A residue with no such atom is not in the universe at all: it is neither scored nor counted against recall. Side chains resolved only in part are therefore scored from what was resolved, without imputation.
  \item[Non-standard residues] A residue type outside the twenty takes the mean of the twenty standard values for each of the seven physicochemical constants, and the neutral prior \(1/2\) for the propensity. Both are neutral in a rank order rather than extreme in it, which is what a sentinel value would have been.
  \item[Chain boundaries] No neighbourhood crosses a chain. Every statistic, every rank and every quantisation is computed inside the chain being scored, which is what allows a chain to be scored without reference to any other structure.
  \item[Heteroatoms] A ligand-leak guard rejects any non-solvent heteroatom before feature extraction, so a receptor that still contains its ligand raises rather than scoring.
\end{description}

\subsection{From the published folds to \NTrainUnits{} and \NUnits{} units}

The dataset is CryptoBench (Skrhak, Novotny, Feidakis, Krivak, Hoksza. CryptoBench: cryptic protein-ligand binding sites dataset and benchmark. Bioinformatics 41(1):btae745 (2025). doi:10.1093/bioinformatics/btae745), taken from its OSF deposit at \url{https://osf.io/pz4a9/}; the SHA-256 of every fetched file is checked against the digest OSF reports. The published split is MMseqs2 clustering at \(10\%\) sequence identity, \(80\%\) coverage, split 80:20. Receptors are fetched from RCSB, scoped to the chain the fold names, and stripped of ligands.

\begin{center}
\begin{tabular}{@{}lrr@{}}
\toprule
 & training & test \\
\midrule
apo structures named by the fold & 885 & 222 \\
apo chains in those structures & 905 & 231 \\
less multi-chain apo chains & $-123$ & $-38$ \\
less chains RCSB would not serve & $-12$ & $-1$ \\
\midrule
evaluated & \NTrainUnits{} & \NUnits{} \\
\bottomrule
\end{tabular}
\end{center}

A unit is one apo chain, so a structure contributing two apo chains contributes two units; that is why the chain counts exceed the structure counts. The multi-chain exclusions are the substantive step. 38 test and 123 training apo chains are compound, named by joining chain letters as in \texttt{M-O-P}. The per-residue metric indexes residues by residue number, which is not unique across chains, so those units cannot be labelled unambiguously without a chain-aware index and are excluded rather than silently mis-scored. The consequence is that the fold reported here is a strict subset of the official test fold and the results do not transfer to multi-chain assemblies. A further one test chain and 12 training chains are absent because RCSB does not serve those entries in PDB format; they are listed by accession in the manifests, the test one being \texttt{7nbc} chain \texttt{CCC}.

The \NTrainUnits{} training units comprise \NTabTrainRes{} residues, of which \(5.759\%\) are labelled cryptic binding residues; that base rate is the value an unaddressed cell returns. The training and test folds share no sequence cluster: the ledger in the compiled artifact records 766 training and 190 test clusters over 770 and 192 units.

